% Lecture example set: SC4001-L01
% Source: 1_foundations pp. 13--14; L1_summary pp. 12--14
% Human verified: Yes

\exercisemeta
  {SC4001-L01-EX}
  {2026-08-14}
  {1\_foundations pp.~13--14；L1\_summary pp.~12--14}
  {答案已根据讲义与讨论核对}
  {已确认（2026-08-16）}

\exercisequestion{SC4001-L01-E01}{Single-neuron Forward Computation}

\textbf{对应知识点：}
\relatedtopic{SC4001-L01-T02}{Artificial Neuron 与 Forward Computation}；
\relatedtopic{SC4001-L01-T03}{Activation Functions 与 Nonlinearity}

\subsubsection*{题目}

一个 neuron（神经元）的参数和 activation function（激活函数）为
\[
  \bm w=
  \begin{bmatrix}2.5\\-0.2\\1.0\end{bmatrix},
  \qquad b=-0.5,
  \qquad f(u)=\frac{0.8}{1+e^{-1.2u}}.
\]
分别求 inputs（输入）
\[
  \bm x^{(1)}=
  \begin{bmatrix}0.8\\2.0\\-0.5\end{bmatrix},
  \qquad
  \bm x^{(2)}=
  \begin{bmatrix}-0.4\\1.5\\1.0\end{bmatrix}
\]
对应的 synaptic input（突触输入）$u$ 和 output（输出）$y$。

\begin{checkedsolution}
对于第一个 input，
\begin{align*}
  u^{(1)}
    &=\bm w^{\mathsf T}\bm x^{(1)}+b \\
    &=2.5(0.8)-0.2(2.0)+1.0(-0.5)-0.5=0.6,\\
  y^{(1)}
    &=\frac{0.8}{1+e^{-1.2(0.6)}}\approx0.538.
\end{align*}
对于第二个 input，
\begin{align*}
  u^{(2)}
    &=2.5(-0.4)-0.2(1.5)+1.0(1.0)-0.5=-0.8,\\
  y^{(2)}
    &=\frac{0.8}{1+e^{-1.2(-0.8)}}\approx0.222.
\end{align*}
\end{checkedsolution}

\textbf{检查点：}先计算 $u$，再代入 $f$；图中连接 constant input（常数输入）$+1$ 的 $-0.5$ 是 bias。Sigmoid-like output（类 Sigmoid 输出）是连续值，$u>0$ 不代表 $y=1$。

\exercisequestion{SC4001-L01-E02}{Discrete Perceptron 与 Odd Parity}

\textbf{对应知识点：}
\relatedtopic{SC4001-L01-T02}{Artificial Neuron 与 Forward Computation}；
\relatedtopic{SC4001-L01-T03}{Activation Functions 与 Nonlinearity}

\subsubsection*{题目}

一个 network（网络）接收 binary input（布尔输入）$\bm x=(x_1,x_2,x_3)\in\{0,1\}^3$。所有 neurons 使用 threshold activation（阈值激活）$\mathbf 1(u>0)$。令
\begin{align*}
  u_1&=x_1+x_2+x_3-\tfrac12,& y_1&=\mathbf 1(u_1>0),\\
  u_2&=x_1+x_2+x_3-\tfrac32,& y_2&=\mathbf 1(u_2>0),\\
  u_3&=x_1+x_2+x_3-\tfrac52,& y_3&=\mathbf 1(u_3>0),\\
  u&=y_1-y_2+y_3-\tfrac12,& y&=\mathbf 1(u>0).
\end{align*}
求该 network 实现的 Boolean function（布尔函数）。

\begin{figure}[htbp]
  \centering
  \resizebox{0.9\textwidth}{!}{\begin{tikzpicture}[
  >=Latex,
  every node/.style={font=\small},
  block/.style={draw=noteBlue,rounded corners,thick,align=center,minimum height=10mm},
  gate/.style={draw=ntuRed,rounded corners,thick,align=center,minimum width=23mm,minimum height=10mm}
]
  \node[block] (input) {$x_1,x_2,x_3$};
  \node[block,right=1.4cm of input] (sum) {$s=x_1+x_2+x_3$};
  \node[gate,right=1.6cm of sum,yshift=1.3cm] (y1) {$y_1=\mathbf 1(s>\tfrac12)$};
  \node[gate,right=1.6cm of sum] (y2) {$y_2=\mathbf 1(s>\tfrac32)$};
  \node[gate,right=1.6cm of sum,yshift=-1.3cm] (y3) {$y_3=\mathbf 1(s>\tfrac52)$};
  \node[block,right=1.8cm of y2] (out) {$y=\mathbf 1\!\left(y_1-y_2+y_3>\tfrac12\right)$};

  \draw[->,thick] (input) -- (sum);
  \draw[->] (sum) -- (y1);
  \draw[->] (sum) -- (y2);
  \draw[->] (sum) -- (y3);
  \draw[->] (y1) -- node[above,sloped] {$+1$} (out);
  \draw[->] (y2) -- node[above] {$-1$} (out);
  \draw[->] (y3) -- node[below,sloped] {$+1$} (out);
\end{tikzpicture}
}
  \caption{讲义 network 的等价紧凑表示：先统计 $s=x_1+x_2+x_3$，再构造三个 threshold features（阈值特征）。}
  \label{fig:sc4001-l01-parity-network}
\end{figure}

\begin{checkedsolution}
令 $s=x_1+x_2+x_3$，即输入中 1 的数量。则 $y_1,y_2,y_3$ 分别表示 $s\ge1$、$s\ge2$、$s=3$。逐个检查 $s\in\{0,1,2,3\}$：
\begin{center}
\begin{tabular}{@{}cccccc@{}}
  \toprule
  $s$ & $y_1$ & $y_2$ & $y_3$ & $u$ & $y$ \\
  \midrule
  0 & 0 & 0 & 0 & $-\tfrac12$ & 0 \\
  1 & 1 & 0 & 0 & $ \tfrac12$ & 1 \\
  2 & 1 & 1 & 0 & $-\tfrac12$ & 0 \\
  3 & 1 & 1 & 1 & $ \tfrac12$ & 1 \\
  \bottomrule
\end{tabular}
\end{center}
因此 network 在输入中有奇数个 1 时输出 1：
\[
  \boxed{y=x_1\oplus x_2\oplus x_3}.
\]
等价的 sum-of-products（乘积项之和）形式为
\[
  y=\bar x_1\bar x_2x_3
   +\bar x_1x_2\bar x_3
   +x_1\bar x_2\bar x_3
   +x_1x_2x_3.
\]
\end{checkedsolution}

\textbf{检查点：}当 $s=2$ 时 $y_1=y_2=1$，两者在 output neuron 中相消，所以 $y=0$。Odd parity（奇校验）不是 linearly separable（线性可分）的；hidden layer（隐藏层）使网络能够实现单个 threshold neuron 无法表示的 Boolean function。
